%
% 14-prinzip.tex -- Die der Methode zugrundeliegenden Prinzipien
%
% (c) 2026 Prof Dr Andreas Müller
%
\subsection{Die der Methode zugrundeliegenden Prinzipien
\label{buch:ratint:bernoulli:subsection:prinzip}}
Es wurde bereits früher festgehalten, dass die dargestellte Methode
für ein Computeralgebrasystem nicht durchführbar ist, weil eine
Methode zur Bestimmung der Faktorisierung des Nenners in Linearfaktoren
oder quadratische reelle Faktoren fehlt.
\index{Computeralgebrasystem}%
Es muss daher eine alternative Vorgehensweise gefunden werden,
die mit einer leichter zu findenden Faktorisierung des Nenners
arbeiten kann.
Die Faktorisierung muss immer noch einige Eigenschaften erfüllen,
damit die einzelnen Schritte verallgemeinert werden können.
Die folgenden Schritt müssen auf die neue Art der Faktorisierung
verallgemeinert werden.

\begin{enumerate}
\item
Die Partialbruchzerlegung wurde dazu verwendet, die 
rationale Funktion in Term mit speziell einfachen Nennern
aufzuspalten.
\index{Partialbruchzerlegung}%
In jedem Term kamen höchstens Potenzen eines einzigen Faktors
des Nenners vor.
Wie in der Darstellung der Partialbruchzerlegung angedeutet
wurde, müssen die Faktoren teilerfremd sein.
\item
Für den Reduktionsschritt
\[
\frac{A}{(z-\alpha)^l}
\quad\to\quad
\frac{\tilde{A}}{(z-\alpha)^{l-1}}
\qquad\text{bzw.}\qquad
\frac{a(z)}{q(z)^l}
\quad\to\quad
\frac{\tilde{a}(z)}{q(z)^{l-1}}
\]
ist eine Rekursionsformel
nötig, die den Exponenten im Nenner reduziert.
\index{Reduktion}%
Wie sich später zeigen wird, beruht sie darauf, dass der
Zähler in der Form
\begin{equation}
p(z) = s(z) \cdot q(z) + t(z)\cdot q'(z)
\label{buch:ratint:printzip:eqn:euklid}
\end{equation}
geschrieben werden kann.
Das aus dem zweiten Summanden resultierende Integral kann
dann durch partielle Integration
\begin{align*}
\int \frac{p(z)}{q(z)^m}\,dz
&=
\int \frac{s(z)}{q(z)^{m-1}}\,dz
+
\int t(z)\cdot \frac{q'(z)}{q(z)^m} \,dz
\\
&=
\int \frac{s(z)}{q(z)^{m-1}}\,dz
+
\int t(z) \frac{d}{dz} \frac{1-m}{q(z)^{m-1}}\,dz
\\
&=
\int \frac{s(z)}{q(z)^{m-1}}\,dz
+
t(z)\frac{1-m}{q(z)^{m-1}}
-
\int \frac{(1-m)t'(z)}{q(z)^{m-1}}\,dz
\end{align*}
vereinfacht werden, was genau ein Rekursionsformel der verlangten
Art ist.
Die Zerlegung \eqref{buch:ratint:printzip:eqn:euklid}
ist genau dann möglich, wenn $q(z)$ und $q'(z)$ teilerfremd
sind, also keine gemeinsame Nullstelle haben.
\index{teilerfremd}%
\item
Es wird ein Algorithmus benötigt, der die Polynom $s(z)$ und
$t(z)$ auf effiziente Art und Weise finden kann.
Um dies im Allgemeinen zu schaffen, müssen wir in
Abschnitt~\ref{buch:ratint:section:euklid} den euklidischen Algorithmus
genauer studieren.% und ihn in Abschnitt
%Abschnitt~\ref{buch:ratint:section:hermite}
%auf quadratfreie Nenner $q(z)$ anwenden.
\item
Es müssen Integrale der Form 
\[
\int \frac{p(z)}{q(z)}\,dz
\]
gefunden werden, wobei $p(z)$ ein Polynom vom Grad $\deg p(z)<\deg q(z)$ ist.
Die Methode von Rothstein-Trager von
Abschnitt~\ref{buch:ratint:section:rothsteintrager}
wird dies möglich machen.
\end{enumerate}

Eine geeignete Faktorisierung, die immer gefunden werden kann,
ist die sogenannte quadratfreie Faktorisierung, die in Abschnitt 
\ref{buch:ratint:section:hermite}
zusammen mit der Reduktion der Integrale von rationalen
Funktionen mit Potenzen des Nennerfaktors vorgestellt wird.


